\documentclass[11pt,a4paper]{article}

% ---------- Page setup ----------
\usepackage[margin=0.7in]{geometry}

% ---------- Fonts & ATS-friendly text extraction ----------
\usepackage[T1]{fontenc}
\usepackage{lmodern}
\usepackage[utf8]{inputenc}
\input{glyphtounicode}
\pdfgentounicode=1

% ---------- Colors & Icons ----------
\usepackage[dvipsnames]{xcolor}
\usepackage{fontawesome5}
\definecolor{accent}{HTML}{1F4E79}
\definecolor{muted}{HTML}{4A4A4A}

% ---------- Links ----------
\usepackage[colorlinks=true,urlcolor=accent,linkcolor=accent]{hyperref}

% ---------- Lists ----------
\usepackage{enumitem}
\setlist[itemize]{leftmargin=*, itemsep=2pt, topsep=3pt, parsep=0pt}

% ---------- Section style settings ----------
\newcommand{\secfontsize}{14}
\newcommand{\seclineheight}{1.2pt}
\newcommand{\secgapabove}{14pt}
\newcommand{\secgaptitle}{6pt}
\newcommand{\secgapbelow}{10pt}
\newcommand{\company}[1]{\textbf{\textcolor{accent}{#1}}}

% ---------- Section style with separator line ----------
\newcommand{\cvsection}[1]{%
    \par\vspace{\secgapabove}%
    {\fontsize{\secfontsize}{\numexpr\secfontsize+3\relax}\selectfont\bfseries\scshape #1}\par
    \vspace{\secgaptitle}%
    {\color{black}\hrule height \seclineheight}%
    \vspace{\secgapbelow}%
    \noindent\ignorespaces
}

\newcommand{\techstack}[1]{{\small\color{muted}#1}}

% ---------- Misc ----------
\pagestyle{empty}
\setlength{\parindent}{0pt}
\raggedbottom

\begin{document}

%====================
% Header
%====================
\noindent
{\fontsize{22}{30}\selectfont\bfseries\color{black} Eric Daniel Hutabarat, S.Kom.}\\[4pt]
{\Large\bfseries\color{muted} Pengembang Fullstack \textcolor{accent}{|} Pengembang Mobile}\\[7pt]
{\small
\textcolor{accent}{\faEnvelope}\hspace{4pt}\href{mailto:ericdaniel.dev@gmail.com}{ericdaniel.dev@gmail.com}
\hspace{16pt}
\textcolor{accent}{\faGithub}\hspace{4pt}\href{https://github.com/erichutabarat}{github.com/erichutabarat}
\hspace{16pt}
\textcolor{accent}{\faLinkedin}\hspace{4pt}\href{https://www.linkedin.com/in/eric-daniel-hutabarat-18192a29a/}{linkedin.com/in/eric-daniel-hutabarat}
}\\[6pt]
{\color{black}\rule{\linewidth}{1.4pt}}

%====================
\cvsection{RINGKASAN PROFESIONAL}
%====================
Software Engineer dengan keahlian di bidang pengembangan web full-stack, rekayasa backend, dan deployment cloud. Mampu mengembangkan REST API berkinerja tinggi dengan Go, membangun antarmuka pengguna yang responsif dengan Next.js, serta men-deploy aplikasi berbasis kontainer menggunakan Docker dan VPS Linux. Memadukan latar belakang akademik yang kuat di bidang Teknik Informatika dengan pengalaman proyek dan magang secara praktis.

%====================
\cvsection{KEAHLIAN TEKNIS}
%====================
\newcommand{\skillitem}[2]{{\bfseries\color{black}#1:} #2\par\vspace{6pt}}

\noindent
\begin{minipage}[t]{0.48\linewidth}\raggedright
\skillitem{Bahasa Pemrograman}{Go, JavaScript, TypeScript, Dart, Python, SQL}
\skillitem{Frontend}{Next.js, ReactJS, Flutter, Tailwind CSS, NiceGUI}
\skillitem{Backend}{Go Gin, Express.js, FastAPI, Prisma ORM, GORM}
\skillitem{Basis Data}{MySQL, PostgreSQL, Firebase, Supabase}
\end{minipage}
\hfill
\begin{minipage}[t]{0.48\linewidth}\raggedright
\skillitem{DevOps}{Docker, Docker Compose, Linux, Nginx, GitHub Actions}
\skillitem{Pengujian}{Jest, Postman, Go testing}
\skillitem{Perangkat}{Git, GitHub, Figma, VS Code, Postman}
\skillitem{IoT}{ESP32, MQTT}
\end{minipage}

%====================
\cvsection{PENGALAMAN PROFESIONAL}
%====================

\textbf{Pengembang Web Full-Stack (Magang)}
\hfill
Juni 2024 -- Agustus 2024 \\
\company{Dinas Perkebunan dan Peternakan Provinsi Sumatera Utara - Medan}
\begin{itemize}
\item Mengembangkan aplikasi web menggunakan Next.js dan Supabase untuk mengumpulkan, mengelola, dan menganalisis tanggapan survei dari para tamu.
\item Merancang dan mengimplementasikan antarmuka pengguna yang responsif serta skema basis data untuk pengumpulan data yang lancar dan pembaruan secara real-time.
\end{itemize}

\vspace{8pt}

\textbf{Pengembang Backend \& Pengembang Aplikasi Mobile}
\hfill
Februari 2025 -- Juni 2026 \\
\company{Biomedis, ITERA - Lampung}
\begin{itemize}
    \item Mengembangkan sistem "Deteksi Cemas", dengan mengimplementasikan aplikasi mobile lintas platform menggunakan Flutter dan Dart untuk interaksi pengguna dan pemantauan secara real-time.
    \item Membangun API backend yang skalabel menggunakan Go dengan Gin dan GORM, serta mengintegrasikan broker Mosquitto MQTT untuk penerimaan data telemetri perangkat keras secara langsung.
    \item Merancang dan mengoptimalkan skema basis data relasional di MySQL untuk mengelola autentikasi pengguna, log data sensor, dan metrik asesmen secara aman.
    \item Mengemas layanan backend, basis data, dan broker MQTT ke dalam kontainer menggunakan Docker dan Docker Compose guna memastikan deployment yang andal dan konsistensi lingkungan.
\end{itemize}

\vspace{8pt}

\textbf{Pengembang GUI Python}
\hfill
Maret 2025 -- Mei 2025 \\
\company{RTK-GPS, Remote}
\begin{itemize}
\item Mengembangkan dasbor pemantauan berbasis web dengan Python menggunakan NiceGUI untuk mesin penyortir hasil panen pertanian yang sudah ada, dengan fokus hanya pada lapisan antarmuka pengguna.
\item Membangun tampilan real-time untuk statistik mesin (jumlah hasil panen, waktu operasi, efisiensi), indikator status komponen (penyortir, komputer, kamera, mesin) dengan deteksi masalah otomatis, serta grafik Plotly interaktif.
\item Mengintegrasikan siaran langsung kamera dari perangkat edge (Jetson) dan menyediakan endpoint REST FastAPI untuk menyajikan data dasbor.
\item Merancang tata letak responsif dengan kelas Tailwind CSS dan menstrukturkan aplikasi dengan modul controller dan layout yang terpisah agar lebih mudah dipelihara.
\end{itemize}

\vspace{8pt}

%====================
\cvsection{PROYEK}
%====================
\textbf{Deteksi Cemas -- Aplikasi Mobile Deteksi Tingkat Kecemasan}\\[2pt]
\techstack{Flutter, Go Gin, MySQL, Docker}
\begin{itemize}
\item Mengembangkan aplikasi mobile untuk mendeteksi tingkat kecemasan menggunakan Flutter, dengan layanan backend yang dibangun menggunakan Golang Gin dan integrasi basis data MySQL.
\item Mengintegrasikan perangkat IoT dan model machine learning yang dikembangkan oleh anggota tim untuk mengumpulkan data fisiologis dan menyediakan analisis tingkat kecemasan.
\item Mengemas layanan backend menggunakan Docker untuk meningkatkan konsistensi deployment dan skalabilitas aplikasi.
\end{itemize}

\vspace{4pt}

\textbf{MySDLC -- Platform Manajemen Siklus Hidup Pengembangan Perangkat Lunak}\\[2pt]
\techstack{Next.js, Go, GORM, MySQL, Docker Compose}
\begin{itemize}
\item Membangun platform manajemen proyek menggunakan Next.js, Go, GORM, dan MySQL.
\item Mengimplementasikan pelacakan proyek, fase SDLC, kolaborasi tim, dan kontrol akses berbasis peran (role-based access control).
\item Men-deploy aplikasi menggunakan Docker Compose dengan build kontainer otomatis.
\end{itemize}

\vspace{4pt}

\textbf{Sistem Pemantauan IoT Automatic Transfer Switch}\\[2pt]
\techstack{ESP32, MQTT, Firebase, PostgreSQL}
\begin{itemize}
\item Mengembangkan sistem pemantauan parameter kelistrikan berbasis ESP32.
\item Mengintegrasikan komunikasi MQTT, Firebase Realtime Database, dan penyimpanan PostgreSQL.
\item Merancang layanan backend untuk pemantauan real-time dan analisis data historis.
\end{itemize}

%====================
\cvsection{PENDIDIKAN}
%====================
\textbf{Institut Teknologi Sumatera}
\hfill
Agustus 2021 -- November 2026 \\
\textit{Sarjana Teknik Informatika} \hfill IPK: 3,32 / 4,00

Berfokus pada Rekayasa Frontend, Pengembangan Backend, Manajemen Basis Data, Pengembangan Aplikasi Mobile, Version Control dengan Git, dan Sistem Internet of Things (IoT).

\vspace{6pt}

\textbf{Dicoding -- Pengembangan Frontend dan Backend}
\hfill
Februari 2024 -- Juli 2024 \\
\textit{Studi Mandiri Bersertifikat} \hfill IPK: 4,00 / 4,00

Menyelesaikan pelatihan intensif di bidang Pengembangan Frontend dan Backend, mencakup pengembangan aplikasi web, pemrograman JavaScript, React, layanan backend, dan praktik pengembangan perangkat lunak.


%====================
\cvsection{SERTIFIKASI}
%====================
\textbf{Belajar Dasar Git dengan GitHub -- Dicoding}
\hfill Desember 2024 \\
Sertifikat: \href{https://www.dicoding.com/certificates/6RPNV0G88Z2M}{Tautan Verifikasi}

\vspace{5pt}

\textbf{Belajar Membuat Front-End Web untuk Pemula -- Dicoding}
\hfill Desember 2024 \\
Sertifikat: \href{https://www.dicoding.com/certificates/1OP8N900QXQK}{Tautan Verifikasi}

\vspace{5pt}

\textbf{Belajar Back-End Pemula dengan JavaScript -- Dicoding}
\hfill Desember 2024 \\
Sertifikat: \href{https://www.dicoding.com/certificates/1OP8WDO28XQK}{Tautan Verifikasi}

\vspace{5pt}

\textbf{Belajar Pengembangan Web Intermediate -- Dicoding}
\hfill Desember 2024 \\
Sertifikat: \href{https://www.dicoding.com/certificates/N9ZOY14GDPG5}{Tautan Verifikasi}

\end{document}